%----------
%	RESULTADOS Y DISCUSIÓN — introducción del capítulo
%----------

Este capítulo evalúa la arquitectura de agentes especializadas con \textit{Agent
Skills} presentada en el Capítulo~\ref{implementation}, aplicada a la
clasificación de las cinco variables de lenguaje sexista del corpus definido en el Capítulo~\ref{sec:methodology}. La
Sección~\ref{sec:iris_experimentacion} describe el diseño experimental y presenta
los resultados sobre las 1.313 noticias con anotación experta, mientras que la
Sección~\ref{sec:discussion} los discute de forma transversal y extrae las
implicaciones para el proyecto IRIS. La Sección~\ref{poc} cierra el capítulo
examinando, a la luz de esos resultados, de qué modo el sistema puede emplearse
en el trabajo editorial y cómo se ha integrado en la herramienta del proyecto.

Este capítulo adopta una óptica concreta para interpretar los resultados. El interés
no recae únicamente en el acuerdo entre el sistema y la anotación, sino en lo que el
acuerdo, y sobre todo el desacuerdo, revela sobre la naturaleza de la tarea. La razón
es que el sexismo lingüístico resulta difícil de percibir, y lo resulta también para
quien está formada para percibirlo, porque la ceguera que produce la cultura patriarcal y el espejismo de
neutralidad androcéntrica actúan sobre usos que la lengua ha normalizado hasta
volverlos poco visibles. De ahí que la variabilidad que se observa entre personas
anotadoras no deba leerse sin más como una divergencia de pareceres sobre qué es
sexismo, sino como el reflejo de esa dificultad. Y de ahí también que un acuerdo bajo
admita dos lecturas que el capítulo se ocupa de separar: la de una tarea cuyo criterio
no llega a estabilizarse y la de un sistema que no detecta parte importante de lo que las expertas sí
señalan de forma mayoritaria, con independencia de quién anote. Es importante señalar que cada una de las piezas con las que trabajamos sólo ha sido marcada por una experta sin que exista, por tanto, la posibilidad de verificar el acuerdo entre ellas. 
